% Original PDF: Section 3, pages 6--13 and the first two lines on page 14.
\section{Results}\label{sec:results}
\begin{theorem}[Convergence under cumulative learning]\label{thm:convergence}
Consider recursion~\eqref{eq:update} under
Assumption~\ref{asp:main}, with either fixed-priority or full-inertia tie-breaking from Section~\ref{sec:algorithm}. Then,
almost surely, $(q_{\pi_k})$ changes only finitely many times and
\[
q_k\longrightarrow q^\star.
\]
Every sufficiently late selected policy $\pi_k$ is optimal. If each state has a unique optimal action,
$\pi_k$ is eventually constant; with optimal ties, constancy of policy labels is not asserted.
\end{theorem}

\subsection{The mean-field argument}
Replace each sampled update by its conditional expectation:
\begin{equation}\label{eq:meanfield}
q_{k+1}(s,a)=q_k(s,a)+\alpha_k\sigma_k(s)\bigl(q_{\pi_k}(s,a)-q_k(s,a)\bigr).
\end{equation}
At a state with $c=\alpha_k\sigma_k(s)>0$, greediness before and after the update gives
\[
\begin{aligned}
q_{\pi_k}(s,\pi_{k+1}(s))-q_{\pi_k}(s,\pi_k(s))
&\ge\frac{1-c}{c}\bigl(q_k(s,\pi_k(s))-q_k(s,\pi_{k+1}(s))\bigr)\\
&\ge0.
\end{aligned}
\]
If $c=0$, the state's table is unchanged. Fixed priorities and full inertia retain its selected
action, so the same policy-improvement inequality holds with equality. Thus Lemma~\ref{lem:improvement} gives
$q_{\pi_{k+1}}\ge q_{\pi_k}$. With finitely many targets, this sequence is eventually a single $\bar q$. Infinite
cumulative learning then makes \eqref{eq:meanfield} converge to $\bar q$. A policy $\widehat\pi$ occurring infinitely often is
greedy for $\bar q=q_{\widehat\pi}$, so Bellman verification gives $\bar q=q^\star$.

Equal action rates are exactly what permits a common coefficient $c$. The sampled
recursion need not improve its target at every change. Instead we will show a uniform
positive conditional probability of improving or stopping during each sufficiently late target
interval.

\subsection{Sampling, filtering, and stability}
\begin{lemma}[From conditional to actual update weight]\label{lem:actual}
Let $J_{k+1}(i)\in\{0,1\}$ be $\mathcal F_{k+1}$-measurable indicators for finitely many coordinates, and put $p_k(i)=\mathbb E[J_{k+1}(i)\mid\mathcal F_k]$. If the
deterministic $\alpha_k\ge0$ satisfy $\sum_k\alpha_k^2<\infty$, then on one probability-one event, for every $i$,
\begin{equation}\label{eq:actual}
\sum_k\alpha_kp_k(i)=\infty\quad\Longrightarrow\quad\sum_k\alpha_kJ_{k+1}(i)=\infty.
\end{equation}
This statement does not require an assumption on the update targets.
\end{lemma}
\begin{proof}
The series $\sum_k\alpha_k(J_{k+1}(i)-p_k(i))$ is a martingale series whose increment variances
sum to at most $\sum_k\alpha_k^2$. It therefore converges almost surely by the $L^2$ martingale convergence
theorem. Subtracting its partial sums from $\sum_{k<n}\alpha_kJ_{k+1}(i)$ proves the implication. Intersect
the probability-one events over the finitely many coordinates.
\end{proof}

For initial-pair indicators, $p_k(i)=\sigma_k(i)$. In particular, a positive initial-pair floor and
$\sum_k\alpha_k=\infty$ give infinite actual update weight. A full first-visit indicator dominates the
initial-pair indicator, so the same conclusion holds for full first visits.

\begin{lemma}[Filtering accumulated errors]\label{lem:filter}
Consider the scalar recursion
\[
x_{k+1}=(1-b_k)x_k+b_ky_k+e_{k+1},\qquad0\le b_k\le1.
\]
For $n>r$, its filtered error satisfies
\begin{equation}\label{eq:filter}
\left|\sum_{k=r}^{n-1}e_{k+1}\prod_{j=k+1}^{n-1}(1-b_j)\right|
\le2\max_{r\le m\le n}\left|\sum_{k=r}^{m-1}e_{k+1}\right|.
\end{equation}
If $\sum_kb_k=\infty$ and $\sum_ke_{k+1}$ converges, eventual bounds $\ell\le y_k\le h$ imply $\ell\le\liminf_kx_k\le
\limsup_kx_k\le h$. In particular, an eventually constant target $y$ gives $x_k\to y$. All statements
are pathwise.
\end{lemma}
\begin{proof}
Unrolling from $r$ gives the initial value with weight $P_{r,n}=\prod_{k=r}^{n-1}(1-b_k)$, the targets
with nonnegative weights summing to $1-P_{r,n}$, and the displayed error. The error weights
belong to $[0,1]$ and are nondecreasing in $k$, so summation by parts gives \eqref{eq:filter}. If $\sum_kb_k=\infty$,
then $P_{r,n}\to0$ for each $r$. Convergence of the error series makes its tail partial sums uniformly
small. First let $n\to\infty$ with $r$ after the target bounds begin, and then let $r\to\infty$.
\end{proof}

\begin{lemma}[Stability and persistent targets]\label{lem:stability}
Under Assumption~\ref{asp:main}, with any greedy tie
selection, there is one probability-one event on which
\[
\sum_k\alpha_k\mathbf1_{\{I_{k+1}=i\}}=\infty\quad(i\in\mathcal I),\qquad
\limsup_k\|q_k\|_\infty\le B.
\]
On the same event, if $q_{\pi_k}=\bar q$ for all sufficiently large $k$, then $q_k\to\bar q=q^\star$.
\end{lemma}
\begin{proof}
For $i=(s,a)$, put
\[
\lambda_{k,i}=\alpha_k\mathbf1_{\{I_{k+1}=i\}},\qquad
\varepsilon_{k+1,i}=\lambda_{k,i}(G_{k+1}-q_{\pi_k}(i)).
\]
Cumulative learning and Lemma~\ref{lem:actual} give $\sum_k\lambda_{k,i}=\infty$ almost surely. The return-moment
assumption makes $\varepsilon_{k+1,i}$ a martingale difference with conditional second moment at most
$c_v\alpha_k^2\sigma_k(s)$. Its series therefore converges almost surely. Intersect these events over the finitely
many coordinates.

On this event the MCES recursion has the form of Lemma~\ref{lem:filter}, with weights $\lambda_{k,i}$, targets
$q_{\pi_k}(i)\in[-B,B]$, and errors $\varepsilon_{k+1,i}$. The lemma gives stability and convergence to any
eventual target $\bar q$. A policy $\widehat\pi$ occurring infinitely often satisfies $q_{\widehat\pi}=\bar q$, and its greedy
inequalities pass to the limit. Lemma~\ref{lem:improvement} gives $\bar q=q^\star$. No conditioning on eventual target
persistence is used.
\end{proof}

\subsection{A barrier on a state's learning clock}
The same sampling factor appears in a gap's drift and variance. Keeping that factor permits
arbitrarily long pauses: when a state is not sampled, both the restoring force and the noise
pause. No divergence of the state's clock is needed for the following local estimate.

\begin{lemma}[State-clock barrier]\label{lem:barrier}
Let $(D_k)$ be a real adapted process with integrable increments, let $r\le\tau$ be stopping times, and let $\sigma_k(s)\in[0,1]$ be $\mathcal F_k$-measurable. The deterministic
steps are in $(0,1]$, square summable, and nonincreasing from $r$ onward. Write
\[
m_k=\mathbb E[D_{k+1}-D_k\mid\mathcal F_k],\qquad
Z_{k+1}=D_{k+1}-D_k-m_k.
\]
Suppose deterministic $\delta,V>0$ satisfy, on $\{r\le k<\tau,D_k\ge0\}$,
\begin{equation}\label{eq:barrierhyp}
m_k\ge\alpha_k\sigma_k(s)(4\delta-D_k),\qquad
\mathbb E[Z_{k+1}^2\mid\mathcal F_k]\le V\alpha_k^2\sigma_k(s).
\end{equation}
Then, on $\{r<\infty,D_r>0\}$,
\begin{equation}\label{eq:barrier}
\mathbb P\bigl(\exists\text{ finite }n\in[r,\tau]:D_n\le0\mid\mathcal F_r\bigr)
\le C_0\frac{\alpha_r}{\min\{D_r,\delta\}}+\frac{C_1}{\delta^2}\sum_{k=r}^{\infty}\alpha_k^2,
\end{equation}
where $C_0=\max\{1,20V(1+(3\delta)^{-1})\}$ and $C_1=16V$. The crossing event includes the update
ending at $\tau$.
\end{lemma}
\begin{proof}
\emph{Before reaching $\delta$.} Suppose $0<D_r<\delta$. If $D_r<\alpha_r$, the first term in \eqref{eq:barrier} is already
at least one. Otherwise define
\[
\nu=\tau\wedge\inf\{n\ge r:D_n\le0\text{ or }D_n\ge\delta\},\qquad
M_n=\sum_{k=r}^{n-1}\mathbf1_{\{k<\nu\}}Z_{k+1},
\]
and the state clock
\begin{equation}\label{eq:clock}
A_s(r,n)=\sum_{k=r}^{n-1}\alpha_k\sigma_k(s).
\end{equation}
A first crossing at $n\le\tau$ before reaching $\delta$ forces
\begin{equation}\label{eq:crossing}
M_n\le-D_r-3\delta A_s(r,n).
\end{equation}
For $b>0$ known at $r$, set $T(b)=\inf\{n\ge r:A_s(r,n)\ge b\}$, possibly infinite. Nonincrease
bounds the overshoot and future variance, at every finite horizon $n\ge r$, by
\[
A_s(r,T(b)\wedge n)\le b+\alpha_r,\qquad
\sum_{k=r}^{T(b)\wedge n-1}\alpha_k^2\sigma_k(s)\le\alpha_r(b+\alpha_r).
\]
Conditional martingale isometry and Doob's $L^2$ inequality therefore give
\begin{equation}\label{eq:doob}
\mathbb P\!\left(\sup_{r\le n\le T(b)}|M_n|\ge z\ \middle|\ \mathcal F_r\right)
\le\frac{4V\alpha_r(b+\alpha_r)}{z^2},\qquad z>0.
\end{equation}
The supremum is over finite indices. First stop at a finite horizon, then increase the horizon;
an unattained supremum is handled by increasing thresholds to $z$ from below. For $\mathcal F_r$-measurable $b,z$, localize on sets where $b,z,z^{-1}$ are bounded. Random $r$ is treated on $\{r=j\}$
and then combined over integers $j$. Thus no finite bound on $T(b)$ or the waiting time has
been assumed.

Partition possible crossing times into $A_s(r,n)<D_r/(3\delta)$ and, for $j\ge1$,
\[
2^{j-1}D_r/(3\delta)\le A_s(r,n)<2^jD_r/(3\delta).
\]
Apply \eqref{eq:doob} with $(b,z)=(D_r/(3\delta),D_r)$ for the first range and $(b,z)=(2^jD_r/(3\delta),2^{j-1}D_r)$
for the others. Because $D_r\ge\alpha_r$, the respective bounds are at most
\[
4V(1+(3\delta)^{-1})\frac{\alpha_r}{D_r},\qquad
16V(1+(3\delta)^{-1})\frac{\alpha_r}{D_r}2^{-j}.
\]
Summing gives the first term in \eqref{eq:barrier}.

\emph{After reaching $\delta$.} Let $\rho=\inf\{n\ge r:D_n\ge\delta\}$, using $\rho=r$ if $D_r\ge\delta$. On $\{\rho<\infty,\rho\le\tau\}$,
restart at $\rho$ and stop at $\eta=\tau\wedge\inf\{n\ge\rho:D_n\le0\}$:
\[
\widetilde M_n=\sum_{k=\rho}^{n-1}\mathbf1_{\{k<\eta\}}Z_{k+1}.
\]
For $D_k\ge\delta$, \eqref{eq:barrierhyp} and $0\le\alpha_k\sigma_k(s)\le1$ imply $D_k+m_k\ge\delta$. If a first subsequent crossing is
at $n$, let $\ell<n$ be the last index with $D_\ell\ge\delta$. All subsequent gaps before $n$ are in $(0,\delta)$ and
have nonnegative mean increments, so
\[
D_n\ge\delta+\widetilde M_n-\widetilde M_\ell.
\]
A crossing forces $\sup_{n\ge\rho}|\widetilde M_n|\ge\delta/2$. Doob's inequality at $\rho$ bounds this probability by
$16V\delta^{-2}\sum_{k\ge\rho}\alpha_k^2$. Multiply by the $\mathcal F_\rho$-measurable event $\{\rho<\infty,\rho\le\tau\}$, average back to $\mathcal F_r$,
and use $\rho\ge r$. This proves the second term. The last-visit index $\ell$ is used only in a pathwise
inequality; we never condition at it. Both parts retain the increment $k=\tau-1$ and therefore
protect the gap through the endpoint itself.
\end{proof}

\subsection{Reference gaps and separately completed seeds}
For a reference policy $\pi$, define
\begin{equation}\label{eq:gaps}
\begin{gathered}
d_{\pi,s,a}=q_\pi(s,\pi(s))-q_\pi(s,a),\qquad
\mathcal B_\pi=\{(s,a):d_{\pi,s,a}>0\},\\
D_k^{\pi,s,a}=q_k(s,\pi(s))-q_k(s,a).
\end{gathered}
\end{equation}
An action in $\mathcal B_\pi$ would worsen the reference action according to $q_\pi$. If these estimated gaps
are all positive, no greedy policy can select a bad action, so policy improvement applies. If
every $\mathcal B_\pi$ is empty, improvement applies between every pair of policies in both directions.
Their state values are equal; choosing a policy taking any prescribed action shows that every
action attains that same value. All policies are optimal with the same target, and Lemma~\ref{lem:stability}
completes the theorem. Otherwise set
\begin{equation}\label{eq:constants}
d_* = \min_{\pi,\ (s,a)\in\mathcal B_\pi}d_{\pi,s,a}>0,\qquad
\delta=d_*/4,\qquad h=g=\delta/2.
\end{equation}

Fix deterministic $K>B$ and restart $r$, and let $\zeta_r=\inf\{k\ge r:\|q_k\|_\infty>K\}$. For a
stopping time $t\ge r$, work on $\{t<\infty,t<\zeta_r\}$, fix $\pi=\pi_t$, and define
\begin{equation}\label{eq:stops}
\tau=\inf\{k>t:q_{\pi_k}\ne q_\pi\},\qquad \chi=\tau\wedge\zeta_r.
\end{equation}
These are stopping times. Return moments give finite second moments for every fixed-time
table by induction in \eqref{eq:update}. For $t\le k<\chi$, the action-value target is exactly $q_\pi$, even if the
selected policy label is different. Equal initial action probabilities give
\begin{align}
\mathbb E[D_{k+1}^{\pi,s,a}-D_k^{\pi,s,a}\mid\mathcal F_k]
&=\alpha_k\sigma_k(s)(d_{\pi,s,a}-D_k^{\pi,s,a}),\label{eq:gapdrift}\\
\operatorname{Var}(D_{k+1}^{\pi,s,a}-D_k^{\pi,s,a}\mid\mathcal F_k)
&\le V\alpha_k^2\sigma_k(s),\qquad V=2\{c_v+(B+K)^2\}.\label{eq:gapvariance}
\end{align}
Only the two coordinates forming the gap can change, their selection events are disjoint,
and each contributes at most $\alpha_k^2\sigma_k(s)\{c_v+(B+K)^2\}$ to the raw second moment. Since
$d_{\pi,s,a}\ge4\delta$, the state-clock barrier applies up to and including $\chi$. All such moment estimates
are used before exit, not conditioned on future containment in the ball.

\paragraph{A favorable local update.}
Suppose the reference action is selected at $s$ and a current
bad gap is below $\delta$. Request one of the following action/return events:
\begin{enumerate}
\item If $q_\pi(s,\pi(s))-q_k(s,\pi(s))\ge d_*/2$, select $\pi(s)$ and require $G_{k+1}\ge q_\pi(s,\pi(s))-h$. The
reference coordinate increases by at least $3\delta\alpha_k/2$.
\item Otherwise select $a$ and require $G_{k+1}\le q_\pi(s,a)+h$. The identity
\[
q_k(s,a)-q_\pi(s,a)=q_k(s,\pi(s))-q_\pi(s,\pi(s))+d_{\pi,s,a}-D_k^{\pi,s,a}>\delta
\]
implies that the bad coordinate decreases by more than $g\alpha_k=\delta\alpha_k/2$.
\end{enumerate}
In either case no other gap against the reference action decreases, and the reference action
remains greedy. For centered $X$ with variance at most $c_v$, the one-sided Chebyshev inequality
gives
\[
\mathbb P(X\ge-h)\ge\frac{h^2}{c_v+h^2},\qquad
\mathbb P(X\le h)\ge\frac{h^2}{c_v+h^2}.
\]
For example, apply Markov's inequality to $(X+c_v/h)^2$ to bound $\mathbb P(X\ge h)$ and repeat for
$-X$; when $c_v=0$, $X=0$ almost surely. Conditional on selecting $s$, its initial action is
uniform, so every requested local update has conditional probability at least
\begin{equation}\label{eq:localprob}
p_0=\frac{h^2}{\max_s|\mathcal A(s)|\,(c_v+h^2)}>0.
\end{equation}
No probability of selecting the state enters this bound.

\begin{lemma}[Separately completed statewise seeds]\label{lem:seeds}
For fixed-priority or full-inertia\linebreak
ties and any positive integer $N$, there is an adapted construction asking for at most $JN$
favorable local updates, starting at $t\ge k_{\mathrm m}$. Along successful requests, every unfinished state
retains its reference action, including at an endpoint reached by a favorable update. If state $s$
completes at time $T_s<\infty$, then
\begin{equation}\label{eq:seedmargin}
D_{T_s}^{\pi,s,a}\ge\min\{\delta,gN\alpha_{T_s}\}\qquad((s,a)\in\mathcal B_\pi).
\end{equation}
Each request, conditional on selecting its unfinished state, has probability at least $p_0$.
\end{lemma}
\begin{proof}
Order the bad actions at each state. Skip an action immediately if its gap is already
at least $\delta$; otherwise process it with at most $N$ favorable updates at that state, stopping
earlier if it reaches $\delta$. Process skips at $t$ and after every favorable update, before the next
sample. A state completes when its list is exhausted. Write $T_s$ for this time, using $T_s=t$ for
a state completed initially and $T_s=\infty$ if the construction stops before completion. These
are stopping times. Every selection of an unfinished state requests a local update; completed
states receive ordinary updates. Stop the construction at the first failed requirement or at $\chi$.

Under full inertia, raising the reference estimate or lowering a competitor retains the
reference action, since it remains greedy. At an unupdated state the table is unchanged, so
the action is also retained. Induction proves the assertion through every unfinished state's
wait.

For fixed priorities, at time $t$ the reference action has highest priority among that state's
maximizers. Leaving the table unchanged, raising its estimate, or lowering a competitor
cannot create a new maximizer ahead of it in the priority order. Therefore that same action
continues to be selected at every unfinished state. Induct over all iterations, including those
updating other states and the final favorable update at completion. This proves the required
retention property directly; a fixed-priority rule is not being identified with full inertia on
arbitrary updates.

Every processed gap either reaches $\delta$ or receives $N$ increases of at least $g\alpha_k$. Later seed
updates at its state cannot lower it. Eventual nonincrease, already in force, gives $\alpha_k\ge\alpha_{T_s}$
for every such update and proves \eqref{eq:seedmargin}. At most $J$ bad actions are processed, so there are
at most $JN$ requests, irrespective of waiting times. Their conditional probability bound is
\eqref{eq:localprob}.
\end{proof}

Completed states can be updated while other states wait. We therefore cannot condition
indiscriminately on all future successful behavior: the completed-state return laws must stay
unchanged. The following bounded change of measure conditions only the seed updates.

\begin{lemma}[Conditioning only the seed updates]\label{lem:measure}
There is an auxiliary law $\widehat{\mathbb P}\ll\mathbb P$
under which all active requirements of Lemma~\ref{lem:seeds} succeed. Conditional state selection is
unchanged. At completed states, and after the construction stops, the conditional action and
return law given the current history and selected state is unchanged. For every future event
$E$, on $\{t<\infty,t<\zeta_r\}$,
\begin{equation}\label{eq:transfer}
\mathbb P(E\mid\mathcal F_t)\ge p_0^{JN}\widehat{\mathbb P}(E\mid\mathcal F_t).
\end{equation}
\end{lemma}
\begin{proof}
Let $S_{k+1}$ be the state component of $I_{k+1}$. At an active unfinished state $s$, let
$p_k(s)$ be the conditional probability of its requirement given $\mathcal F_k,S_{k+1}=s$. Then $p_k(s)\ge p_0$.
If the state has zero sampling probability, set $p_k(s)=p_0$; this value is never used on a
positive-probability draw.

Start $L_t=1$. At a requested update of $s$, multiply $L$ by the indicator of success divided
by $p_k(s)$; at other iterations multiply it by one. Every factor has conditional mean one given
$\mathcal F_k,S_{k+1}$. Hence $(L_k)$ is a martingale. At most $JN$ factors differ from one, even though their
times need not be bounded, so $0\le L_k\le p_0^{-JN}$. Bounded convergence gives $\mathbb E[L_\infty\mid\mathcal F_t]=1$.
Define $d\widehat{\mathbb P}=L_\infty\,d\mathbb P$. For random $t$, use likelihood one before $t$, on $\{t=\infty\}$, and off $\{t<\zeta_r\}$,
and combine the constructions on $\{t=j\}$ over integers $j$. The law agrees with $\mathbb P$ on $\mathcal F_t$
where the construction is active.

A failure makes $L_\infty=0$, so active requirements succeed almost surely under the auxiliary
law. To check conditional laws, use $\mathbb E[L_\infty\mid\mathcal F_{k+1}]=L_{k+1}$; future factors do not alter
the one-step likelihood. Its current factor has conditional mean one given the selected
state, preserving state selection. At a completed state the current factor is identically one,
so its conditional action and return law is unchanged as well. The statement holds on
auxiliary-almost-every history, where $L_k>0$. Absolute continuity preserves the almost-sure
algorithmic identities, including the chosen fixed-priority or full-inertia tie rule. Finally,
\[
\widehat{\mathbb P}(E\mid\mathcal F_t)=\mathbb E[L_\infty\mathbf1_E\mid\mathcal F_t]
\le p_0^{-JN}\mathbb P(E\mid\mathcal F_t),
\]
which proves the bound. This auxiliary law is used only in the proof; the algorithm still
samples its original returns.
\end{proof}

\subsection{A stopped positive probability of progress}
Let $W(\pi)=\sum_{i\in\mathcal I}q_\pi(i)$. Its possible values form a finite ordered set. A distinct componentwise larger target has strictly larger $W$, although different targets may have equal
$W$.

\begin{lemma}[Stopped improvement bound]\label{lem:progress}
Under the assumptions of Theorem~\ref{thm:convergence}, for
every fixed $K>B$ there are deterministic $k_0$ and $p>0$ such that for every deterministic
restart $r$ and stopping time $t\ge\max\{r,k_0\}$, on $\{t<\infty,t<\zeta_r\}$,
\begin{equation}\label{eq:progress}
\mathbb P\!\left(
\begin{aligned}
&\{\zeta_r\le\tau,\zeta_r<\infty\}\text{ or }\{\tau=\zeta_r=\infty\}\text{ or}\\
&\{\tau<\zeta_r,W(\pi_\tau)>W(\pi_t)\}
\end{aligned}
\ \middle|\ \mathcal F_t\right)\ge p.
\end{equation}
Here $\tau$ is the next target change in \eqref{eq:stops}.
\end{lemma}
\begin{proof}
Fix $V$ from \eqref{eq:gapvariance} and $C_0,C_1$ from the state-clock barrier. Choose a positive integer
$N$ with $2JC_0/(\delta N)\le1/4$. Then choose deterministic $k_0\ge k_{\mathrm m}$ so that, for every $t\ge k_0$,
\begin{equation}\label{eq:choices}
N\alpha_t\le2,\qquad\frac{JC_1}{\delta^2}\sum_{k=t}^{\infty}\alpha_k^2\le\frac14.
\end{equation}
The radius and variance constants are fixed before choosing $N$.

Use the auxiliary law of Lemma~\ref{lem:measure}. For each bad action whose state completes at
$T_s<\chi$, apply the barrier at $T_s$, stopped at $\chi$. At that completed state the original sampling
and return laws are preserved, so \eqref{eq:gapdrift}--\eqref{eq:gapvariance} hold under $\widehat{\mathbb P}$. No such assertion is made for
a conditioned return at an unfinished state. Nonincrease and \eqref{eq:seedmargin}--\eqref{eq:choices} imply that its
conditional crossing probability is at most
\begin{equation}\label{eq:crossprob}
\frac{2C_0}{\delta N}+\frac{C_1}{\delta^2}\sum_{k=t}^{\infty}\alpha_k^2.
\end{equation}
Indeed $gN\alpha_{T_s}\le\delta$, so $\alpha_{T_s}/\min\{D_{T_s}^{\pi,s,a},\delta\}\le1/(gN)$. Multiply the conditional estimate by
the $\mathcal F_{T_s}$-measurable event $\{T_s<\chi\}$ and average back to $\mathcal F_t$. A union bound over at most
$J$ pairs gives probability at most $1/2$ of any completed gap failing through $\chi$. This needs
neither independence nor simultaneous completion. If $T_s=\chi<\infty$, its final favorable update
already retains the reference action, so no subsequent barrier is needed.

On the complementary event, either exit occurs first, the target persists without exit,
or $\tau<\zeta_r$ is finite. In the third case, states unfinished just before the endpoint retain the
reference action, including those completed at the endpoint. Previously completed states
cannot choose a bad action. Thus
\[
q_\pi(s,\pi_\tau(s))\ge q_\pi(s,\pi(s))\quad\text{for every }s.
\]
Policy improvement gives $q_{\pi_\tau}\ge q_\pi$ with distinct targets, and hence a strict increase of $W$.
There is no need to prove that every seed completes: at a finite endpoint unfinished states
are safe, and an infinite endpoint without exit is already a successful outcome. Transfer back
by \eqref{eq:transfer}, obtaining \eqref{eq:progress} with $p=\frac12p_0^{JN}$.

If $\mathcal B_\pi$ is empty, the gap conditions and barrier are vacuous and policy improvement
holds for every next selected policy. The same lower bound remains valid. An exit \emph{after} an
unsuccessful finite target change is not counted as success for that earlier interval.
\end{proof}

\subsection{Finitely many target changes and convergence}
\begin{proof}[Proof of Theorem~\ref{thm:convergence}]
Fix $K>B$ and constants $k_0,p$ from Lemma~\ref{lem:progress}. For a deterministic
$r\ge k_0$, start at $t_0=r$ and absorb immediately if $\|q_r\|_\infty>K$. Enumerate successive target
changes
\[
t_{n+1}=\inf\{k>t_n:q_{\pi_k}\ne q_{\pi_{t_n}}\}
\]
only while $t_n<\zeta_r$. All finite continuing times are stopping times. Stop counting if exit
occurs before or at the next target change, or if the target persists forever without exit.

At every alive interval, absorption or a strict increase in $W$ has conditional probability
at least $p$. By induction and the tower property, absorption within the next $m$ intervals
or $m$ consecutive strict increases has conditional probability at least $p^m$. On an absorbed
branch the continuation probability is one; on a finite successful-transition branch condition
at its stopping time and apply the remaining bound. This uses no independence and never
conditions at an infinite time on a continuing branch.

Let $L$ be the number of distinct values of $W$. There cannot be $L$ consecutive strict
increases. Absorption within the next $L$ intervals therefore has conditional probability at least
$p^L$. Applying the bound at successive surviving blocks gives probability at most $(1-p^L)^j$
of surviving $j$ such blocks. In particular the probability of infinitely many target changes
before exit is zero.

Intersect these probability-one conclusions over all deterministic integers $r\ge k_0$. Lemma~\ref{lem:stability}
says that almost every path remains inside the ball after some such integer. On that path,
choose $r$ with $\zeta_r=\infty$; the original sequence consequently has only finitely many target
changes. The integer is selected pathwise only after the countable intersection. We do not
condition on an eventual containment time or claim that it is a stopping time. Nor is an
unconditional geometric tail for the original process's total number of changes asserted.

Lemma~\ref{lem:stability} now gives $q_k\to q^\star$. Every nonoptimal action has a strictly positive optimality
gap. Finiteness of the action sets and convergence exclude all such actions from sufficiently
late greedy policies. Those policies are greedy for $q^\star$ and hence optimal by Bellman verification.
If every state has a unique optimal action, their labels are eventually constant.
\end{proof}
